\section{Comparative Register of Textual and Interpretive Positions}

The register below maps every directly inspected biblical text, version, Second Temple retelling, Jewish and Catholic commentary, and New Testament echo examined across Genesis 15:1--21.

\begin{fourstable}{Source \& Date}{Passage}{Hour, Vision, and Going Outside}{Bearing on Daylight Stars}
Masoretic Hebrew Text (Kittel 1909) & Gen 15:1, 5, 12, 17--18 & Vision (\latin{ba-ma\d{h}azeh}, v.~1); \latin{wayyiqtol} chain in v.~5 with no elapsed-time measure; sun about to set (\latin{l\={a}v\={o}'}, v.~12); \latin{b\={a}'\={a}h} and dark mist (\latin{'al\={a}\d{t}\={a}h}, v.~17); \latin{ba-yy\={o}m ha-h\={u}'} (v.~18). & Permits night-to-next-evening sequence or visionary scope; does not require daylight at v.~5. \\
\addlinespace[0.3em]
Septuagint (Swete 1909, Alexandrinus; Brenton) & Gen 15:1, 5, 11--12, 17--18 & Vision (\latin{en horamati}); Abram sits by carcasses (v.~11); sunset trance (\latin{ekstasis}, v.~12); sun still setting with flame (\latin{phlox}), no dark mist (v.~17); Alexandrinus reads ``there'' (\latin{ekei}, v.~18), Sixtine/Brenton ``in that day.'' & Weakens sunset-to-darkness contrast and pulls v.~12 toward prophetic ecstasy; raises no daylight tension. \\
\addlinespace[0.3em]
Clementine Vulgate \& Douay-Rheims (Jerome; Challoner) & Gen 15:1, 5, 11--12, 17--18 & Vision (\latin{per visionem}); Abram drives birds away (v.~11); imperfect \latin{occumberet} and \latin{sopor} (v.~12); pluperfect \latin{occubuisset sol, facta est caligo tenebrosa} (v.~17); \latin{in illo die} (v.~18). & Sharpens afternoon-to-nightfall tenses between vv.~12 and 17, yet Jerome never flags v.~5 as daytime. \\
\addlinespace[0.3em]
Syriac Peshitta & Gen 15:5, 12, 17--18 & Preserves bringing Abram outside and counting stars (v.~5), sun approaching setting (vv.~12, 17), and ``on that day'' (v.~18) with no inserted interval. & Mirrors the Hebrew sequence without adding any daytime-stars gloss. \\
\addlinespace[0.3em]
Book of Jubilees (2nd c.~BC) & Jub 14:1--20 & Explicitly writes Abram \emph{beheld the stars} (14:5); dates sacrifice to mid-third month (14:10); sunset at 14:13, 17; retains ``on that day'' (14:18). & Assumes visible stars at 14:5 without a day-break or daylight paradox. \\
\addlinespace[0.3em]
Flavius Josephus (AD~37--c.~100) & \emph{Antiquities} I.10.3 & Keeps the star-multitude promise and sacrifice; omits going outside, looking up, sunset, and darkness. & Neutral abridgment; supplies no support for daylight stars. \\
\addlinespace[0.3em]
Isaac Abarbanel (1437--1508) & Gen 15, prefatory QQ.~7 \& 16 (Warsaw 1862) & Weighs earlier visionary commentators (\latin{be-mar'eh ha-nevu'ah}) and \emph{Shabbat} 156a (``brought him out from his astrology''); frames the waking dilemma (Q.~7) and questions why v.~12 names the hour despite the Ran's ``lamp at midday'' argument (Q.~16). Body resolution not contained in inspected excerpt. & Earliest explicit statement of the daylight-stars dilemma; rejects waking daytime star-counting outright as impossible (\latin{'i 'efshar}). \\
\addlinespace[0.3em]
St. Clement of Rome (c.~AD~96) & \emph{1 Clement} 10 & Quotes Gen 15:5--6 from the Septuagint beside Gen 12:1--3 and 13:14--16 as a monument of Abraham's faith. & Earliest Christian citation; raises no chronological tension. \\
\addlinespace[0.3em]
St. Ambrose (c.~340--397) & \emph{De excessu} II.95 & Alludes to Gen 15:5: ``To them Wisdom showed the heaven illuminated with stars.'' & Casually assumes a starlit sky; \emph{De Abraham} not directly collated. \\
\addlinespace[0.3em]
St. Jerome (c.~347--420) & \emph{Quaestiones Hebr.} ch.~15 & Scrutinizes Heb.~15:2--3, 7, 10--11, 12 (\latin{thardema}/\latin{sopor}), and 16 word by word; passes over v.~5 in silence. & Decisive philological silence from the Vulgate's translator. \\
\addlinespace[0.3em]
St. Augustine (354--430) & \emph{City of God} X.8, XVI.23--24; \emph{Gen. ad litt.} XII.6 & Declares ``all these things were said and done in a vision'' (\latin{Haec omnia in visu}); stars signify celestial felicity and refute Aratus and Eudoxus; sunset = end of the world. & Whole-chapter visionary umbrella eliminates any terrestrial daylight clash. \\
\addlinespace[0.3em]
St. Bede, Alcuin, \& Angelomus & Commentaries on Gen 15 & Bede reads ``indoors'' as fleshly Israel and ``outside'' (\latin{foras}) as the universal Church shining as stars; Alcuin asks six questions on Gen 15, none on v.~5. & Unbroken early medieval Latin silence on any daytime-stars difficulty. \\
\addlinespace[0.3em]
St. Thomas Aquinas (1225--1274) & \emph{ST} II-II q.~173 a.~3, q.~174 a.~6; III q.~30 a.~3 ad~3; \emph{In Gal.} 3:16 & Distinguishes sensible, imaginary, and intellectual prophecy; places Abraham at the head of pre-Mosaic revelation (q.~174 a.~6); cites Gen 15:12's \latin{horror} as imaginary vision disturbing bodily powers. & Provides the scholastic framework for prophetic vision; no daylight-stars issue. \\
\addlinespace[0.3em]
Cornelius a Lapide (1567--1637) & \emph{In Pentateuchum} (1700), Gen 15 & Deduces at v.~5: ``Therefore it was then night---cloudless, clear, and starry'' (citing the \latin{dioptra}); explains v.~12 sleep by ``excessive daytime labor'' (\latin{ex nimio labore diurno}). & Implicitly spans night (v.~5) through the next day's labor (vv.~9--11) to sunset (v.~12); denies daylight at v.~5. \\
\addlinespace[0.3em]
New Testament Echoes & Rom 4:3, 18--22; Gal 3:6--9; Heb 11:11--12 & Paul quotes ``So shall thy seed be'' and v.~6, naming Abraham's dead body and Sarah's dead womb (so Chrysostom, Hom.~9 on Rom); Heb 11:12 blends Gen 15:5 and 22:17. & Complete apostolic silence on invisible daytime stars; faith confronts sterility with God's Word. \\
\end{fourstable}
